% Chinese questions and concise solutions. Compile using XeLaTeX.
\documentclass[UTF8,10pt,letterpaper,fontset=fandol]{ctexart}
\usepackage[margin=1in]{geometry}
\usepackage{amsmath,amssymb}
\usepackage{enumitem,needspace}
\usepackage[unicode,hidelinks]{hyperref}
\pagestyle{plain}
\raggedbottom
\setlength{\parindent}{0pt}
\setlength{\parskip}{1pt}
\setlength{\emergencystretch}{2em}
\linespread{1.1}
\AtBeginDocument{%
\setlength{\abovedisplayskip}{6pt}%
\setlength{\belowdisplayskip}{6pt}%
\setlength{\abovedisplayshortskip}{3pt}%
\setlength{\belowdisplayshortskip}{3pt}}
\setcounter{secnumdepth}{0}
\setlist[enumerate]{label=(\arabic*),leftmargin=2em,itemsep=0.25em,topsep=0.2em}
\ctexset{section={format=\Large\bfseries,beforeskip=1.2ex,afterskip=0.8ex}}
\newcommand{\Solution}{\par\Needspace{4\baselineskip}\textbf{解答}\par}
\begin{document}
\begin{center}
{\LARGE 作业 6}\\[0.6em]
2026 年 10 月 23 日\\[0.25em]
姓名：YuanHao Tang\quad 学号：3120265621
\end{center}

\section{第 1 题}
设 $\mathcal A$ 是四维线性空间 $\mathcal V$ 到自身的线性变换，$\alpha_1,\alpha_2,\alpha_3,\alpha_4$ 是 $\mathcal V$ 的一组基。
$\mathcal A$ 关于此基的矩阵表示为
\[
A=\begin{pmatrix}2&0&2&1\\6&1&4&4\\10&0&0&4\\7&0&-7&2\end{pmatrix}.
\]
求：
\begin{enumerate}
\item 矩阵 $A$ 的特征值与特征向量；
\item 线性变换 $\mathcal A$ 的特征值及相应特征向量。
\end{enumerate}
\Solution
\begin{enumerate}
\item 特征多项式为 $\det(\lambda I-A)=(\lambda-1)^3(\lambda-2)$。又
$\operatorname{rank}(A-I)=\operatorname{rank}(A-2I)=3$，对应特征值 $1,2$ 的全部特征向量分别为
\[
\boxed{x_1=c(0,1,0,0)^T,\qquad x_2=c(1,2,1,-2)^T,\qquad c\ne0.}
\]
\item 特征值同为 $1,2$。将坐标还原为原空间中的向量，全部相应特征向量为
\[
\boxed{v_1=c\alpha_2,\qquad v_2=c(\alpha_1+2\alpha_2+\alpha_3-2\alpha_4),\qquad c\ne0.}
\]
\end{enumerate}

\section{第 2 题}
设
\[
A=\begin{pmatrix}4&6&0\\-3&-5&0\\-3&-6&1\end{pmatrix}.
\]
求 $A$ 的不变多项式、自然标准形 $N$，以及复数情形和实数情形的 Jordan 标准形。
\Solution
对特征矩阵 $\lambda I-A$，有 $D_1=1$。因 $\operatorname{rank}(A-I)=1$，全部二阶子式均含 $\lambda-1$，其中一个为 $-6(\lambda-1)$，故 $D_2=\lambda-1$。
又 $D_3=\det(\lambda I-A)=(\lambda-1)^2(\lambda+2)$，所以不变多项式为
\[
\boxed{d_1=1,\qquad d_2=\lambda-1,\qquad d_3=(\lambda-1)(\lambda+2)}.
\]
去掉单位项，对非常数不变多项式取列伴随块，得到
\[
\boxed{N=\begin{pmatrix}1&0&0\\0&0&2\\0&1&-1\end{pmatrix}}.
\]
特征值 $1$ 的几何重数为 $2$，$-2$ 为单根，因此两种 Jordan 形均为
\[
\boxed{J_{\mathrm{complex}}=J_{\mathrm{real}}=\operatorname{diag}(1,1,-2)}.
\]

\newpage
\section{第 3 题}
设矩阵 $A$ 的初等因子为
\[
\begin{gathered}
(\lambda+1),\quad(\lambda+1)^2,\quad(\lambda+1+j),\quad(\lambda+1-j),\\
(\lambda+1)^2,\quad(\lambda+1+j)^2,\quad(\lambda+1-j)^2.
\end{gathered}
\]
求：
\begin{enumerate}
\item 特征矩阵 $\lambda E-A$ 的 Smith 标准形 $S$；
\item $A$ 的自然标准形 $N$、复 Jordan 标准形 $J_{\mathrm{complex}}$ 及实 Jordan 标准形 $J_{\mathrm{real}}$。
\end{enumerate}
\Solution
令 $j^2=-1$、$x=\lambda+1$、$q=x^2+1$。初等因子总次数为 $11$。
\begin{enumerate}
\item 按各根的指数升序排列，在左侧补零至三项：
\[
\begin{array}{c|ccc}
\text{根}&1&2&3\\\hline
-1&1&2&2\\
-1-j&0&1&2\\
-1+j&0&1&2
\end{array}
\]
故非常数不变多项式为 $\psi_1=x$、$\psi_2=x^2q$、$\psi_3=x^2q^2$，于是
\[
\boxed{S=\operatorname{diag}(I_8,\psi_1,\psi_2,\psi_3)}.
\]
\item 用 $C(p)$ 表示首一多项式 $p=\lambda^m+c_{m-1}\lambda^{m-1}+\cdots+c_0$ 的列伴随矩阵：次对角线为 $1$，最后一列为 $-(c_0,\ldots,c_{m-1})^T$，其余为零。则
\[
\boxed{N=\operatorname{diag}\bigl(-1,C(\psi_2),C(\psi_3)\bigr)},
\]
其中
\[
\begin{aligned}
\psi_2&=\lambda^4+4\lambda^3+7\lambda^2+6\lambda+2,\\
\psi_3&=\lambda^6+6\lambda^5+17\lambda^4+28\lambda^3+28\lambda^2+16\lambda+4.
\end{aligned}
\]
定义
\[
J_2(\rho)=\begin{pmatrix}\rho&1\\0&\rho\end{pmatrix},\qquad
K_\pm=\operatorname{diag}\bigl(-1\pm j,J_2(-1\pm j)\bigr).
\]
每个初等因子对应一个 Jordan 块，所以
\[
\boxed{J_{\mathrm{complex}}=\operatorname{diag}\bigl(-1,J_2(-1),J_2(-1),K_-,K_+\bigr)}.
\]
实数情形将共轭块配对。令
\[
R=\begin{pmatrix}-1&1\\-1&-1\end{pmatrix},\qquad
R_2=\begin{pmatrix}R&I_2\\0&R\end{pmatrix},
\]
得到
\[
\boxed{J_{\mathrm{real}}=\operatorname{diag}\bigl(-1,J_2(-1),J_2(-1),R,R_2\bigr)}.
\]
以上各标准形均为 $11\times11$，Jordan 块的排列次序可改变。
\end{enumerate}
\end{document}
